\section{Median-of-Means (MoM) Estimator}
The following sections are adapted from \cite{devroye2016sub} and \cite{lugosi2019mean}


\subsection{Sub-Gaussian Estimators}
\begin{definition}{Sub-Gaussian Estimators}
Let $(X_i)_{i=1}^{n}$ be a random sample with $\E[X_i] = \mu$ and $\Var[X_i] = \sigma^2$, and let $\bar{\mu}(X)$ denote an estimator for $\mu$. For some constant $L > 0$, $\bar{\mu}(X)$ is said to be a sub-Gaussian estimator for $\mu$ if 
\begin{equation}
\label{eq:subgaussian}
    \Prob[\abs{\bar{\mu}(X) - \mu} > L\sigma \sqrt{\frac{\ln(1/\delta)}{n}}] \leq \delta \quad \text{for all } \delta \in (0,1).
\end{equation}
\end{definition}

\subsection{The Empirical Mean Estimator}
Let $(X_i)_{i=1}^{n}$ be a random sample with $\E[X_i] = \mu$ and $\Var[X_i] = \sigma^2 < \infty$. 
The empirical mean is defined as
\[
\bar{\mu} = \frac{1}{n}\sum_{i=1}^{n} X_i.
\]
By the Central Limit Theorem,
\[
\sqrt{n}(\bar{\mu} - \mu) \overset{d}{\to} \mathcal{N}(0, \sigma^2),
\]
which implies that for any fixed $\delta \in (0,1)$,
\[
\Prob[\abs{\bar{\mu} - \mu} > \frac{\sigma \Phi^{-1}(1-\delta/2)}{\sqrt{n}}] \to \delta 
\quad \text{as } n \to \infty.
\]
Using the bound $\Phi^{-1}(1-\delta/2) \leq \sqrt{2\ln(2/\delta)}$, it follows that 
\[
\lim_{n \to \infty} \Prob[\abs{\bar{\mu} - \mu} > \sigma\sqrt{\frac{2\ln(2/\delta)}{n}}] \leq \delta,
\]
so $\bar{\mu}$ is an asymptotically sub-Gaussian estimator in the sense of \eqref{eq:subgaussian} with constant $L = \sqrt{2}$.

While $\bar{\mu}$ is asymptotically sub-Gaussian, it is not a sub-Gaussian estimator for finite samples; Chebyshev's inequality gives 
\[
\Prob[\abs{\bar{\mu} - \mu} > \sigma\sqrt{\frac{1}{n\delta}}] \leq \delta
\]
shows that the dependence on $\delta$ is polynomial rather than logarithmic, 
which is substantially worse than \eqref{eq:subgaussian} for small $\delta$. As such, the empirical mean is insufficient as a sub-Gaussian estimator for random variables with only finite first and second moments.

\subsection{The Median-of-Means Estimator}
The median-of-means (MoM) estimator is not new, being proposed in various forms, first by \cite{nemirovskij1983problem} and \cite{jerrum1986random}. The algorithm works as follows

\begin{algorithm}[H]
\caption{Median-of-Means Estimator}
\KwIn{Sample $(X_i)_{i=1}^{n}$, number of blocks $k$}
\KwOut{Estimate $\hat{\mu}$}

Partition $(X_i)_{i=1}^{n}$ into $k$ blocks $B_1, \dots, B_k$ of size $m = \lfloor n/k \rfloor$\;

\For{$j = 1, \dots, k$}{
    Compute the block mean $\bar{\mu}_j = \frac{1}{m}\sum_{i \in B_j} X_i$\;
}

\Return{$\hat{\mu} = \med(\bar{\mu}_1, \dots, \bar{\mu}_k)$}\;
\end{algorithm}

The MoM estimator satisfies
\[
    \Prob[\abs{\hat{\mu} - \mu} > \sigma\sqrt{\frac{4}{m}}] \leq e^{-k/8},
\]
and setting $k = \lceil 8\ln(1/\delta) \rceil$, this becomes
\begin{equation}
    \Prob[\abs{\hat{\mu} - \mu} > \sigma\sqrt{\frac{32\ln(1/\delta)}{n}}] \leq \delta 
    \quad \text{for all } \delta \in (0,1)
\end{equation}
which is sub-Gaussian with $L = \sqrt{32}$